% Ad M. Brutum 2.5 (ad-brutum-02-05)
% Source: https://raw.githubusercontent.com/PerseusDL/canonical-latinLit/master/data/phi0474/phi059/phi0474.phi059.perseus-lat1.xml
% Fetched by scripts/fetch_latin.py

% Dateline (Perseus): Scr. Romae xviii K. Mai. a. 711 (43).

\ciceroLetterOpener{Cicero Bruto salutem}

\ciceroSection{1} quae litterae tuo nomine recitatae sint id . April. in senatu eodemque tempore Antoni credo ad te scripsisse tuos; quorum ego nemini concedo. sed nihil necesse erat eadem omnis, illud necesse me ad te scribere quid sentirem tota de constitutione huius belli et quo iudicio essem quaque sententia. voluntas mea, Brute, de summa re publica semper eadem fuit quae tua, ratio quibusdam in rebus (non enim omnibus) paulo fortasse vehementior. scis mihi semper placuisse non rege solum sed regno liberari rem publicam; tu lenius immortali omnino cum tua laude; sed quid melius fuerit magno dolore sensimus, magno periculo sentimus. recenti illo tempore tu omnia ad pacem quae oratione confici non poterat, ego omnia ad libertatem quae sine pace nulla est. pacem ipsam bello atque armis effici posse arbitrabar. studia non deerant arma poscentium; quorum repressimus impetum ardoremque restinximus.

\ciceroSection{2} itaque res in eum locum venerat ut, nisi Caesari Octaviano deus quidam illam mentem dedisset, in potestatem perditissimi hominis et turpissimi M. Antoni veniendum fuerit, quocum vides hoc tempore ipso quod sit quantumque certamen. id profecto nullum esset, nisi tum conservatus esset Antonius. sed haec omitto; res enim a te gesta memorabilis et paene caelestis repellit omnis reprehensiones, quippe quae ne laude quidem satis idonea adfici possit. exstitisti nuper vultu severo; exercitum, copias, legiones idoneas per te brevi tempore comparasti. di immortales! qui ille nuntius, quae illae litterae, quae laetitia senatus, quae alacritas civitatis erat! nihil umquam vidi tam omnium consensione laudatum. erat exspectatio reliquiarum Antoni, quem equitatu legionibusque magna ex parte spoliaras. ea quoque habuit exitum optabilem. nam tuac litterae quae recitatae in senatu sunt et imperatoris et militum virtutem et industriam tuorum, in quibus Ciceronis mei, declarant. quod si tuis placuisset de his litteris referri et nisi in tempus turbuientissimum post discessum Pansae consulis incidissent, honos quoque iustus et debitus dis immortalibus decretus esset.

\ciceroSection{3} ecce tibi Idib. April. advolat mane Celer Pilius, qui vir, di boni, quam gravis, quam constans, quam bonarum in re publica partium! hic epistulas adfert duas, unam tuo nomine, alteram Antoni; dat Servilio tribuno plebis, ille Cornuto. recitantur in senatu. ANTONIVS PROCOS. Magna admiratio, ut si esset recitatum DOLABELLA IMPERATOR ; a quo quidem venerant tabellarii, sed nemo Pili similis qui proferre litteras auderet aut magistratibus reddere. tuae recitantur breves illae quidem sed in Antonium admodum lenes. vehementer admiratus senatus. mihi autem non erat explicatum quid agerem. falsas dicerem? quid si tu eas adprobasses?

\ciceroSection{4} confirmarem ? non erat dignitatis tuae. itaque ille dies silentio. postridie autem cum sermo increbruisset, Piliusque oculos vehementius hominum offendisset, natum omnino est principium a me. de proconsule Antonio multa. Sestius causae non defuit post me, cum quanto suum filium, quanto meum in periculo futurum diceret, si contra proconsulem arma tulissent. Nosti hominem; causae non defuit. dixerunt etiam alii. Labeo vero noster nec signum tuum in epistula nec diem adpositum nec te scripsisse ad tuos, ut soleres. hoc cogere volebat falsas litteras esse et, si quaeris, probabat.

\ciceroSection{5} nunc tuum est consilium, Brute, de toto genere belli. video te lenitate delectari et eum putare fructum esse maximum praeclare quidem, sed aliis rebus, aliis temporibus locus esse solet debetque clementiae. nunc quid agitur, Brute? templis deorum immortalium imminet hominum egentium et perditorum spes, nec quicquam aliud decernitur hoc bello nisi utrum simus necne. cui parcimus aut quid agimus? hic ergo consulimus quibus victoribus vestigium nostrum nullum relinquetur? nam quid interest inter Dolabellam et quemvis Antoniorum trium? quorum si cui parcimus, duri fuimus in Dolabella. haec ut ita sentiret senatus populusque Romanus, etsi res ipsa cogebat, tamen maxima ex parte nostro consilio atque auctoritate perfectum est. tu si hanc rationem non probas, tuam sententiam defendam, non relinquam meam. neque dissolutum a te quicquam homines exspectant nec crudele. huius rei moderatio facilis est. ut in duces vehemens sis. in milites liberalis.

\ciceroSection{6} Ciceronem meum, mi Brute, velim quam plurimum tecum habeas. virtutis disciplinam meliorem reperiet nullam quam contemplationem atque imitationem tui. XVIII Kalend. Maias.
